\scanpage{54}
% archive sheet number: 53
\[
\begin{aligned}
A_{[p,q]}&=\varprojlim A_{[p]}\otimes_{A_p}A_q/q^n A_q\\
&=\varprojlim A_{[p]}/q^n A_{[\unsure{p}]}\otimes_{A_p}A_q/q^n A_q
\end{aligned}
\]
\[A_{[p,q,\unsure{r|v|\varpi}]}\qquad\text{etc}\]

$X$ \uline{ensemble ordonné}

$S_X$ ensemble simplicial associé, dont les $n$-simplexes sont les chaines
\[
\begin{aligned}
x_0\leq x_1\leq\cdots\leq x_n\\
\longleftarrow
\end{aligned}
\]
[i.e. les applications croissantes de $[0,n]=\underline\Delta_n$ dans $X$]

À tout $\sigma\in S_X$ est associé un $A_\sigma$

À tout \unsure{opérateur|opération} \deleted{face} \inserted{\unsure{simplicial}} $u:\underline\Delta_m\longrightarrow\underline\Delta_n$ (\uline{croissant}) est associé un hom
\diagram{\begin{tikzcd}[ampersand replacement=\&]
\deleted{A_\sigma}\&S_X^n\arrow[r,"u^*"]\&S_X^m
\end{tikzcd}}

et si $\sigma\in S_X^{\unsure{n|m}}$, d'où $u^*(\sigma)\in S_X^m$, on a un \unsure{morphisme}
\[
\varphi(\sigma,u):\quad A_\sigma\longleftarrow A_{u^*\sigma},
\]
De façon \unsure{fonctorielle}.
